% ============================================================
%  第三卷 数学附录（§a–§f）
%  底本：高教社中文版第6版（严肃 译，喀兴林 校），原书 pp.590–604
%  本文件以 \chapter*{数学附录} 开头；节号为字母 §a–§f，
%  公式编号为 (a.1) 式；文件末尾恢复节号/公式编号定义（附录为
%  全书最后一个文件，其后无正文）。
%  插图：附录图1–6 在 figures_manifest.md 中无条目（manifest 仅收
%  正文图1–51），按任务约定不插图；正文"附录图N"引用照排。
%  正文圈码 9 处（§a×2、§b×1、§c 标题×1、§d×3、§e×1、§f×1）已由脚注代理
%  全部转录为 \footnote{}（按 pages300 页面实况）：p603×1、p604×1、p606×1、
%  p608×1、p610×1、p611×2、p613×1、p617×1。§c 标题圈码按原书排在节标题后
%  （\protect\footnote，目录用可选参数中的无圈码标题）；p611 两脚注分别系于
%  (d.7) 前"熟知公式"与 (d.8) 前"回到原来的零值"两处。公式类脚注：
%  p606 含 Ai(x)=Φ(x)/√π 式；p613 含 F(α,γ,z)=lim_{β→∞}F(α,β,γ,z/β) 式。
%  照录备查：p603 脚注"vol. II，；"逗号分号连用系原书排印如此。
%  已修正的 OCR/原书错误（均对照 原书转换/pages/p-603…p-617.png 核实）：
%   1. §a V=e^{xt}QZ：OCR "e^{*}QZ" 按原书 p590 恢复；同页 P,Q,Z 式
%      中 i^2→t^2、x_t→xt（OCR 的 t/i 混淆）。
%   2. §a "基中所取的积分路线C"："基中"系原书 p590 排印错误，径改
%      "其中"。
%   3. §a 曲线两端"(t=+∞ 或 l=−∞)"：l 系原书 p591 排印错误，径改
%      t=−∞；OCR 的 C_ι/C′_∥（下标1被认成希腊字母ι、符号∥）一律
%      恢复 C_1、C′_1、C′_2；行22802 "沿 C′_1 和 C′ 的两个积分"
%      第二个漏下标 2，按 p591 补为 C′_2。
%   4. (a.4) OCR "（Ω−1)Ω"→(−1)^n。
%   5. (b.2) OCR 整式乱码（"2x/13·2x/3xx×∫"），按原书 p592 重建为
%      y(x)=常数×∫_C e^{xt−t^3/3}dt。
%   6. (b.5) OCR "φ((x))"→Φ(x)；"φ(−1.02)=0.95"→Φ(−1.02)=0.95
%      （艾里函数记号统一为 \varPhi）。
%   7. (b.6) OCR "x>0 Hf, φ∈b(x)" 按 p594 恢复 "x>0 时，Φ(x)"；
%      I、K 之和后的"≡"号 OCR 误作"="。
%   8. I_ν(x)=i^{−ν}J_ν(ix)：OCR 指数负号丢失，按 p594 恢复 i^{−ν}。
%   9. (c.1)–(c.5) OCR 下标 ι→l；d^nP_l/(dcosθ)^m→d^m（m 次导）。
%  10. (c.12) OCR dσ_1→do_1、δ_{l′}→δ_{ll′}（按 p596）。
%  11. Y_lm 表：OCR e^{εiφ}→e^{±iφ}，上标乱码（\scriptsize ±2iφ 等）
%      恢复为 e^{±2iφ}、e^{±3iφ}。
%  12. §d 函数定义行 OCR "0(t)=t(t−1)"→Q(t)；围线积分各处 e^{\imath}
%      /e^{\prime} 均为 e^{tz}、e^t。
%  13. (d.13) OCR 行首"＝＝"重复按原书两行推导排（=结转下行）。
%  14. (d.16) OCR 分母 "1↓↓z"→1!z（按 p599）。
%  15. (d.17) OCR "(ℓ−1)^{′*1}"→(−1)^{s+1}；节末句补句点。
%  16. §e "β^{(D}"→β①（脚注圈码）；(e.4) 段 "u,"→u_1；(e.3) 段
%      OCR 窜入的 C′（"确实满足 C′ 方程"）按 p601 删；"可得下列
%      关系式 22" 的 "22" 为 OCR 页码噪声，删。
%  17. (e.11) 编号乱码 "(e_n^a+S)^a"→(e.11)，公式本体按 p602 重建。
%  18. (f.1) 推导行：OCR "e^{−cν·a·k·x}"→e^{−(λ−kt)z}、λ^{⋯1}→
%      λ^{−ν−1}、(1−kt/λ)^{γ−1}→^{−ν−1}（按 p602）。
%  19. (f.5) 编号 OCR 漏字母 f；"−(m−1)l"→−(m−1)!、
%      "nαi^{−n−n−1}(m−1)⋯(m−1)"→n!(m−n−1)⋯(m−1)。
%  20. (f.6) 推导行 e^{−(λ−k)z}→e^{−(λ−kt)z}、d^α→d^n、"I["→"["、
%      "(¬^{−1})^n"→(−1)^n（按 p603）。
%  21. (f.11) F 的参数 OCR 缺 "−n,α′,"，按 p604 补全。
%  22. (f.16) "Γ_γ^{θ}"→J_γ^{s0}、"\bar{]}_γ^{s−1,0}(α,α′+1)"→
%      J_γ^{s−1,0}(α,α′+1)；末段 "J_γ^{∘p}"→J_γ^{s0}。
%  23. 替换箭头：OCR 丢失的 → 按原书补全（t→t+z、tz→t、t→1−t/z、
%      t→z−t、t→t/(1−z+zt)、t→λt/(k′t+λ−k′)）；"代人"→"代入"。
%  24. 半角标点 , ( ) : 按文体约定改全角；正文裸排的 x、n、t、s、p
%      等一律收入 $...$；数学环境内保持 ASCII。
%  遗留问题：
%   - (f.7) 求和项分子 "n(n−1)⋯(n−s)" 照原书 p603 直排（俄文原版
%     同），未按物理改写；
%   - (b.7) 后 "x=0 时." 的句点照原书版面；
%   - 记号：正体函数 F/P/Q/Z/L/H/G/J（含 §f 积分 J）按 OCR 与成品
%     章统一排 \mathrm{}；I、K、Z_{1/3} 与 §a 的 P、Q、Z、V 按原书
%     为斜体；Y_{lm} 斜体，与正文 (28.7)（ch04）一致。
% ============================================================
\chapter*{数学附录}
\addcontentsline{toc}{chapter}{数学附录}
\setcounter{section}{0}
\renewcommand{\thesection}{\alph{section}}
\renewcommand{\theequation}{\alph{section}.\arabic{equation}}
% ctex 的节号选项存于内部 \CTEX@thesection（标题/目录用它而非
% \thesection），须一并改为字母，标题才排成 "§a 厄米多项式"
\ctexset{section={number=\alph{section}}}

\section{厄米多项式}

方程式
\begin{equation}
  y''-2xy'+2ny=0
\end{equation}
属于用拉普拉斯方法\footnote{可参考，E.~Goursat, Cours d'Analyse Math\'ematique, Gauthier--Villars, Paris, vol.~II，；或斯米尔诺夫著《高等数学教程》（中译本：第三卷第三分册，第五章~§107），高等教育出版社.}能解的类型.

这个方法可用于以下类型的任一线性方程：
\begin{equation*}
  \sum_{m=0}^{n}\left(a_m+b_mx\right)\frac{\mathrm{d}^my}{\mathrm{d}x^m}=0,
\end{equation*}
式中的系数不超过 $x$ 的一次幂，其求解步骤如下.我们作下列多项式
\begin{equation*}
  P(t)=\sum_{m=0}^{n}a_mt^m,\qquad Q(t)=\sum_{m=0}^{n}b_mt^m
\end{equation*}
及由它们所组成的函数
\begin{equation*}
  Z(t)=\frac{1}{Q}\exp\int\frac{P}{Q}\,\mathrm{d}t,
\end{equation*}
它确定到一个常数因子为止.所考虑方程之解可表成下列复积分形式：
\begin{equation*}
  y=\int_C Z(t)\mathrm{e}^{xt}\,\mathrm{d}t,
\end{equation*}
其中所取的积分路线 $C$ 使积分为有限且不等于零，并且，当 $t$ 走完曲线 $C$ 以后（曲线 $C$ 可以是封闭的，也可以是不封闭的），函数
\begin{equation*}
  V=\mathrm{e}^{xt}QZ
\end{equation*}
回到原值.

在方程式(a.1)的情形下，我们有
\begin{equation*}
  P=t^2+2n,\quad Q=-2t,\quad
  Z=-\frac{1}{2t^{n+1}}\mathrm{e}^{-\frac{t^2}{4}},\quad
  V=\frac{1}{t^n}\mathrm{e}^{xt-\frac{t^2}{4}},
\end{equation*}
故其解为
\begin{equation}
  y=\int\exp\left(xt-\frac{t^2}{4}\right)\frac{\mathrm{d}t}{t^{n+1}}.
\end{equation}

在物理应用中我们只需考虑 $n>-\frac{1}{2}$ 的情况.对这些 $n$ 值而言，积分路线可取曲线 $C_1$ 或 $C_2$（附录图1）；这些曲线满足所需条件，因为函数 $V$ 在曲线两端（$t=+\infty$ 或 $t=-\infty$）等于零\footnote{这些路线对负整数的 $n$ 不适用，因积分(a.2)沿这些路线将恒等于零.}.

我们来求参量 $n$ 的值，使得方程(a.1)之解对 $x$ 的所有有限值而言是有限的.当 $x\to\pm\infty$ 时，此解不比 $x$ 的有限次幂更快地趋向无穷大.我们先考虑非整数的 $n$ 值.此时(a.2)式沿 $C_1$ 和 $C_2$ 的积分给出方程(a.1)的两个独立解.按照 $t=2(x-u)$ 的关系引进变量 $u$ 后，沿 $C_1$ 的积分可变换成以下形式，式中略去了一个常数因子，
\begin{equation}
  y=\mathrm{e}^{x^2}\int_{C'_1}\frac{\mathrm{e}^{-u^2}}{(u-x)^{n+1}}\,\mathrm{d}u,
\end{equation}
其中积分是沿 $u$ 复平面上的 $C'_1$ 曲线进行的，如附录图2所示.

当 $x\to+\infty$ 时，整个积分路线 $C'_1$ 移向无穷远处，(a.3)式中的积分按 $\mathrm{e}^{-x^2}$ 规律趋于零.但当 $x\to-\infty$ 时，积分路线延及整个实轴，(a.3)式中的积分并不指数式地趋于零，从而 $y(x)$ 函数按 $\mathrm{e}^{x^2}$ 规律趋向无穷大.同理，容易证明(a.2)沿曲线 $C'_2$ 的积分当 $x\to+\infty$ 时指数式地发散.

当 $n$ 为正整数时（包括零），沿积分路线的直线部分的积分相互抵消，而(a.3)式中沿 $C'_1$ 和 $C'_2$ 的两个积分现在变成一个绕 $u=x$ 点的封闭曲线的积分，故解为
\begin{equation*}
  y(x)=\mathrm{e}^{x^2}\oint\frac{\mathrm{e}^{-u^2}}{(u-x)^{n+1}}\,\mathrm{d}u,
\end{equation*}
这个解满足所述条件.按照熟知的求解析函数导数的柯西公式
\begin{equation*}
  f^{(n)}(x)=\frac{n!}{2\pi i}\oint\frac{f(t)}{(t-x)^{n+1}}\,\mathrm{d}t,
\end{equation*}
可知 $y(x)$ 除了一个常数因子外就是一个厄米多项式
\begin{equation}
  \mathrm{H}_n(x)=(-1)^n\mathrm{e}^{x^2}\frac{\mathrm{d}^n}{\mathrm{d}x^n}\mathrm{e}^{-x^2}.
\end{equation}

$\mathrm{H}_n$ 多项式按 $x$ 的降幂展开，呈下列形式：
\begin{equation}
  \mathrm{H}_n(x)=(2x)^n-\frac{n(n-1)}{1}(2x)^{n-2}
  +\frac{n(n-1)(n-2)(n-3)}{1\times2}(2x)^{n-4}-\cdots,
\end{equation}
它只含宇称与 $n$ 相同的各种 $x$ 幂次.我们写出前面几个厄米多项式：
\begin{equation}
\begin{aligned}
  &\mathrm{H}_0=1,\quad \mathrm{H}_1=2x,\quad
  \mathrm{H}_2=4x^2-2,\quad \mathrm{H}_3=8x^3-12x,\\
  &\mathrm{H}_4=16x^4-48x^2+12.
\end{aligned}
\end{equation}

计算归一化积分时，我们把(a.4)代入 $\mathrm{e}^{-x^2}\mathrm{H}_n$，并进行分部积分 $n$ 次，得
\begin{equation*}
  \int_{-\infty}^{\infty}\mathrm{e}^{-x^2}\mathrm{H}_n^2(x)\,\mathrm{d}x
  =\int_{-\infty}^{\infty}(-1)^n\mathrm{H}_n(x)
  \frac{\mathrm{d}^n}{\mathrm{d}x^n}\mathrm{e}^{-x^2}\,\mathrm{d}x
  =\int_{-\infty}^{\infty}\mathrm{e}^{-x^2}
  \frac{\mathrm{d}^n}{\mathrm{d}x^n}\mathrm{H}_n\,\mathrm{d}x.
\end{equation*}
但是 $\mathrm{d}^n\mathrm{H}_n/\mathrm{d}x^n$ 是一个常数 $2^nn!$. 故
\begin{equation}
  \int_{-\infty}^{\infty}\mathrm{e}^{-x^2}\mathrm{H}_n^2(x)\,\mathrm{d}x=2^nn!\sqrt{\pi}.
\end{equation}

\section{艾里函数}

方程
\begin{equation}
  y''-xy=0
\end{equation}
也属于拉普拉斯类型（见§a）.根据一般方法，我们作下列函数：
\begin{equation*}
  P=t^2,\quad Q=-1,\quad
  Z=-\mathrm{e}^{-\frac{1}{3}t^3},\quad
  V=\mathrm{e}^{xt-\frac{t^3}{3}},
\end{equation*}
故其解可表成下列形式：
\begin{equation}
  y(x)=\text{常数}\times\int_C\mathrm{e}^{xt-\frac{1}{3}t^3}\,\mathrm{d}t,
\end{equation}
所选的积分路线 $C$，必须使 $V$ 函数在该曲线的两端都等于零.这两个端点因而必须在 $t$ 复平面上的 $\mathrm{Re}\,t^3>0$ 区域内（附录图3中的阴影区）趋向无穷远.

对所有 $x$ 而言均为有限的解，可取图中所示的积分曲线 $C$ 得到.这条曲线可以任意移动，只要它的两个端点仍留在原来的两个阴影区（附录图3中的 I 和 III）内并趋向无穷远.我们注意到，如果所取的曲线位于 III 和 II 两个阴影区内，当 $x\to\infty$ 时，所得之解为无穷大.

把 $C$ 移到虚轴的位置，(b.2)函数呈下列形式（用 $t=iu$ 代入）：
\begin{equation}
  \varPhi(x)=\frac{1}{\sqrt{\pi}}\int_0^{\infty}
  \cos\left(ux+\frac{1}{3}u^3\right)\,\mathrm{d}u.
\end{equation}
函数(b.2)中的常数已取作 $-i/2\sqrt{\pi}$，这样得到的 $\varPhi(x)$ 函数称为艾里函数\footnote{我们按照 В.~А.~福克（Фок）所设的定义；见 Г.~Я.~Яковлева, 艾里函数表. 莫斯科：科学出版社 1969. $\varPhi(x)$ 是福克所定义的两个函数之一，记作 $V(x)$. 文献中还会遇到艾里函数的另一种定义，它与(b.3)差一个常因子：
\begin{equation*}
  Ai(x)=\frac{1}{\sqrt{\pi}}\,\varPhi(x).
\end{equation*}}.

用鞍点法计算积分(b.2)，可得 $x$ 值很大时的 $\varPhi(x)$ 渐近式.对 $x>0$ 而言，被积函数中的指数幂当 $t=\pm\sqrt{x}$ 时有一个极值，这个被积函数的“最陡下降方向”平行于虚轴.因此欲求 $x$ 为很大正值时的渐近式，我们把指数幂展成 $t+\sqrt{x}$ 的幂级数，并沿平行于虚轴的 $C_1$ 直线（图3）积分；$C_1$ 和虚轴的间距 $OA=\sqrt{x}$. 用 $t=-\sqrt{x}+iu$ 代入，我们得
\begin{equation*}
  \varPhi(x)\approx\frac{1}{2\sqrt{\pi}}\int_{-\infty}^{+\infty}
  \exp\left(-\frac{2}{3}x^{3/2}-\sqrt{x}\,u^2\right)\mathrm{d}u,
\end{equation*}
故
\begin{equation}
  \varPhi(x)\approx\frac{1}{2x^{1/4}}\exp\left(-\frac{2}{3}x^{3/2}\right).
\end{equation}
可见 $x$ 为很大正值时，$\varPhi(x)$ 函数指数式递减.

为了求 $x$ 为很大负值时的渐近式，我们注意到，$x<0$ 时指数幂在 $t=i\sqrt{|x|}$ 和 $t=-i\sqrt{|x|}$ 处有一极值，这两点处的最陡下降方向与实轴分别夹 $-\frac{1}{4}\pi$ 角和 $\frac{1}{4}\pi$ 角.取折线 $C_3$（距离为 $OB=\sqrt{|x|}$）为积分路线，经过某些简单变换后，我们得
\begin{equation}
  \varPhi(x)=\frac{1}{|x|^{1/4}}\sin\left(\frac{2}{3}|x|^{3/2}+\frac{\pi}{4}\right).
\end{equation}
因此在 $x$ 为很大负值的区域内，$\varPhi(x)$ 函数具有振荡性.可以指出，$\varPhi(x)$ 函数的第一个极大值（最大极值）为 $\varPhi(-1.02)=0.95$.

艾里函数可以用 $\frac{1}{3}$ 阶的贝塞尔函数来表述.很容易证实(b.1)方程具有下列解：
\begin{equation*}
  \sqrt{x}\,Z_{1/3}\left(\frac{2}{3}x^{3/2}\right),
\end{equation*}
其中 $Z_{1/3}(x)$ 为 $\frac{1}{3}$ 阶贝塞尔方程的任意解.与(b.3)式相同的解为
\begin{equation}
\begin{aligned}
  &x>0\ \text{时},\quad
  \varPhi(x)=\frac{\sqrt{\pi x}}{3}\left[I_{-1/3}\left(\frac{2}{3}x^{3/2}\right)
  -I_{1/3}\left(\frac{2}{3}x^{3/2}\right)\right]
  \equiv\sqrt{\frac{x}{3\pi}}\,K_{1/3}\left(\frac{2}{3}x^{3/2}\right)\\
  &x<0\ \text{时},\quad
  \varPhi(x)=\frac{\sqrt{\pi|x|}}{3}\left[\mathrm{J}_{-1/3}\left(\frac{2}{3}|x|^{3/2}\right)
  +\mathrm{J}_{1/3}\left(\frac{2}{3}|x|^{3/2}\right)\right]
\end{aligned}
\end{equation}
其中
\begin{equation*}
  I_\nu(x)=i^{-\nu}\mathrm{J}_\nu(ix),\qquad
  K_\nu(x)=\frac{\pi}{2\sin\nu\pi}\left[I_{-\nu}(x)-I_\nu(x)\right],
\end{equation*}
利用递推公式
\begin{equation*}
\begin{aligned}
  &K_{\nu-1}(x)-K_{\nu+1}(x)=-\frac{2\nu}{x}K_\nu(x),\\
  &2K'_\nu(x)=-K_{\nu-1}(x)-K_{\nu+1}(x).
\end{aligned}
\end{equation*}
很易求出艾里函数的导数
\begin{equation}
  \varPhi'(x)=-\frac{x}{\sqrt{3\pi}}K_{2/3}\left(\frac{2}{3}x^{3/2}\right),\quad(x>0)
\end{equation}
$x=0$ 时，
\begin{equation*}
\begin{gathered}
  \varPhi(0)=\frac{\sqrt{\pi}}{3^{2/3}\Gamma\left(\frac{2}{3}\right)}=0.629,\\
  \varPhi'(0)=-\frac{3^{1/6}\Gamma\left(\frac{2}{3}\right)}{2\sqrt{\pi}}=-0.459.
\end{gathered}
\end{equation*}
附录图4绘的是艾里函数的曲线图形.

\section[勒让德多项式]{勒让德多项式\protect\footnote{球谐函数的理论已有许多论述极佳的数学文献.为参考计，我们在这里仅给出一些主要的关系式，不准备系统讨论这种函数的理论.}}

勒让德多项式 $\mathrm{P}_l(\cos\theta)$ 由下式定义：
\begin{equation}
  \mathrm{P}_l(\cos\theta)=\frac{1}{2^ll!}
  \frac{\mathrm{d}^l}{(\mathrm{d}\cos\theta)^l}
  \left(\cos^2\theta-1\right)^l
\end{equation}
满足下列微分方程
\begin{equation}
  \frac{1}{\sin\theta}\frac{\mathrm{d}}{\mathrm{d}\theta}
  \left(\sin\theta\frac{\mathrm{dP}_l}{\mathrm{d}\theta}\right)
  +l(l+1)\mathrm{P}_l=0.
\end{equation}

连带勒让德多项式由下式定义：
\begin{equation}
  \mathrm{P}_l^m(\cos\theta)=\sin^m\theta
  \frac{\mathrm{d}^m\mathrm{P}_l(\cos\theta)}{(\mathrm{d}\cos\theta)^m}
  =\frac{1}{2^ll!}\sin^m\theta
  \frac{\mathrm{d}^{l-m}}{(\mathrm{d}\cos\theta)^{l-m}}
  \left(\cos^2\theta-1\right)^l,
\end{equation}
或者，等价地定义为
\begin{equation}
  \mathrm{P}_l^m(\cos\theta)=(-1)^m
  \frac{(l+m)!}{(l-m)!\,2^ll!}\sin^{-m}\theta
  \frac{\mathrm{d}^{l-m}}{(\mathrm{d}\cos\theta)^{l-m}}
  \left(\cos^2\theta-1\right)^l.
\end{equation}
其中 $m=0,1,\cdots,l$，连带多项式满足下列方程
\begin{equation}
  \frac{1}{\sin\theta}\frac{\mathrm{d}}{\mathrm{d}\theta}
  \left(\sin\theta\frac{\mathrm{dP}_l^m}{\mathrm{d}\theta}\right)
  +\left[l(l+1)-\frac{m^2}{\sin^2\theta}\right]\mathrm{P}_l^m=0.
\end{equation}

勒让德多项式的归一化积分式为
\begin{equation*}
  \int_{-1}^{1}\left[\mathrm{P}_l(\mu)\right]^2\mathrm{d}\mu
\end{equation*}
（$\mu=\cos\theta$）.把(c.1)代入上式，并进行 $l$ 次分部积分后，得
\begin{equation*}
  \frac{(-1)^l}{2^{2l}(l!)^2}\int_{-1}^{1}\left(\mu^2-1\right)^l
  \frac{\mathrm{d}^{2l}}{\mathrm{d}\mu^{2l}}\left(\mu^2-1\right)^l\mathrm{d}\mu
  =\frac{(2l)!}{2^{2l}(l!)^2}\int_{-1}^{1}\left(1-\mu^2\right)^l\mathrm{d}\mu.
\end{equation*}
作变换 $u=\frac{1}{2}(1-\mu)$，使该积分变成欧拉 $B$ 积分，结果为
\begin{equation}
  \int_{-1}^{1}\left[\mathrm{P}_l(\mu)\right]^2\mathrm{d}\mu
  =\frac{2}{2l+1}.
\end{equation}
同理，容易证明 $l$ 值不同的 $\mathrm{P}_l(\mu)$ 函数是彼此正交的.
\begin{equation}
  \int_{-1}^{1}\mathrm{P}_l(\mu)\mathrm{P}_{l'}(\mu)\,\mathrm{d}\mu=0,
  \qquad l\neq l'
\end{equation}

连带勒让德多项式的归一化积分可用类似方法算出.把 $\left[\mathrm{P}_l^m(\mu)\right]^2$ 写成(c.3)和(c.4)式的乘积，作 $l-m$ 次分部积分；结果得
\begin{equation}
  \int_{-1}^{1}\left[\mathrm{P}_l^m(\mu)\right]^2\mathrm{d}\mu
  =\frac{2}{2l+1}\frac{(l+m)!}{(l-m)!}.
\end{equation}
容易证明 $l$ 不同（$m$ 相同）的 $\mathrm{P}_l^m$ 函数彼此正交：
\begin{equation}
  \int_{-1}^{1}\mathrm{P}_l^m(\mu)\mathrm{P}_{l'}^m(\mu)\,\mathrm{d}\mu=0,
  \qquad l\neq l'
\end{equation}

三个勒让德函数的乘积的积分计算，见§107.

下列加法定理对勒让德多项式成立.设 $\gamma$ 为球角 $\theta,\varphi$ 和 $\theta',\varphi'$ 所决定的两个空间方向的夹角：$\cos\gamma=\cos\theta\cos\theta'+\sin\theta\sin\theta'\cos(\varphi-\varphi')$，则有
\begin{equation}
\begin{aligned}
  \mathrm{P}_l(\cos\gamma)={}&\mathrm{P}_l(\cos\theta)\mathrm{P}_l(\cos\theta')\\
  &+\sum_{m=1}^{l}2\frac{(l-m)!}{(l+m)!}
  \mathrm{P}_l^m(\cos\theta)\mathrm{P}_l^m(\cos\theta')
  \cos m(\varphi-\varphi'),
\end{aligned}
\end{equation}
这个定理也可用(28.7)定义的球谐函数表成
\begin{equation}
  \mathrm{P}_l(\boldsymbol{n}\cdot\boldsymbol{n}')
  =\frac{4\pi}{2l+1}\sum_{m=-l}^{l}
  Y_{lm}^{*}(\boldsymbol{n}')Y_{lm}(\boldsymbol{n}).
\end{equation}
$\boldsymbol{n}$ 和 $\boldsymbol{n}'$ 是两个单位矢量.$Y_{lm}(\boldsymbol{n})$ 的宗量是 $\boldsymbol{n}$ 相对于某一固定坐标系的球角.

如果(c.10)式乘以 $\mathrm{P}_l(\cos\theta)$ 并对 $\mathrm{d}o=\sin\theta\,\mathrm{d}\theta\,\mathrm{d}\varphi$ 积分.式右所有含有 $\cos m(\varphi-\varphi')$ 因子的各项对 $\varphi$ 积分后等于零；应用(c.6)和(c.7)，我们得
\begin{equation*}
  \int\mathrm{P}_l(\cos\gamma)\mathrm{P}_{l'}(\cos\theta)\,\mathrm{d}o
  =\delta_{ll'}\frac{4\pi}{2l+1}\mathrm{P}_l(\cos\theta').
\end{equation*}
此式可写成下列对称形式，
\begin{equation}
  \int\mathrm{P}_l(\boldsymbol{n}_1\cdot\boldsymbol{n}_2)
  \mathrm{P}_{l'}(\boldsymbol{n}_1\cdot\boldsymbol{n}_3)\,\mathrm{d}o_1
  =\delta_{ll'}\frac{4\pi}{2l+1}
  \mathrm{P}_l(\boldsymbol{n}_2\cdot\boldsymbol{n}_3),
\end{equation}
其中 $\boldsymbol{n}_1,\boldsymbol{n}_2,\boldsymbol{n}_3$ 是三个单位矢量，积分是对 $\boldsymbol{n}_1$ 的方向而言的.

最后，我们给出前几个归一化的 $Y_{lm}$ 球函数：
\begin{equation*}
\begin{aligned}
  Y_{00}&=\frac{1}{\sqrt{4\pi}},\qquad
  &Y_{1,\pm1}&=\mp i\sqrt{\frac{3}{8\pi}}\sin\theta\,\mathrm{e}^{\pm i\varphi},\\
  Y_{10}&=i\sqrt{\frac{3}{4\pi}}\cos\theta,\qquad
  &Y_{2,\pm1}&=\pm\sqrt{\frac{15}{8\pi}}\cos\theta\sin\theta\,\mathrm{e}^{\pm i\varphi},\\
  Y_{20}&=\sqrt{\frac{5}{16\pi}}\left(1-3\cos^2\theta\right),\qquad
  &Y_{2,\pm2}&=-\sqrt{\frac{15}{32\pi}}\sin^2\theta\,\mathrm{e}^{\pm 2i\varphi},\\
  Y_{30}&=-i\sqrt{\frac{7}{16\pi}}\cos\theta\left(5\cos^2\theta-3\right),\\
  Y_{3,\pm1}&=\pm i\sqrt{\frac{21}{64\pi}}\sin\theta\left(5\cos^2\theta-1\right)\mathrm{e}^{\pm i\varphi},\\
  Y_{3,\pm2}&=-i\sqrt{\frac{105}{32\pi}}\cos\theta\sin^2\theta\,\mathrm{e}^{\pm 2i\varphi},\qquad
  &Y_{3,\pm3}&=\pm i\sqrt{\frac{35}{64\pi}}\sin^3\theta\,\mathrm{e}^{\pm 3i\varphi}.
\end{aligned}
\end{equation*}

\section{合流超几何函数}

合流超几何函数定义为级数
\begin{equation}
  \mathrm{F}(\alpha,\gamma,z)=1+\frac{\alpha}{\gamma}\frac{z}{1!}
  +\frac{\alpha(\alpha+1)}{\gamma(\gamma+1)}\frac{z^2}{2!}+\cdots,
\end{equation}
此级数对 $z$ 的所有有限值都是收敛的；参量 $\alpha$ 是任意的，而参量 $\gamma$ 假定不等于零或负整数.如果 $\alpha$ 是一个负整数（或零），则 $\mathrm{F}(\alpha,\gamma,z)$ 变成一个 $|\alpha|$ 次多项式.

$\mathrm{F}(\alpha,\gamma,z)$ 函数满足微分方程
\begin{equation}
  zu''+(\gamma-z)u'-\alpha u=0,
\end{equation}
它很易直接验证\footnote{$\gamma$ 为负整数时的方程(d.2)不需专门研究，因为通过变换［此变换给出(d.3)式］它可以化成 $\gamma$ 为正整数的情形.}.用 $u=z^{1-\gamma}u_1$ 代入方程，上式可变换成另一个同样类型的方程.
\begin{equation}
  zu_1''+(2-\gamma-z)u_1'-(\alpha-\gamma+1)u_1=0.
\end{equation}
由此可知 $\gamma$ 为非整数时，方程(d.2)尚有特解 $z^{1-\gamma}\mathrm{F}(\alpha-\gamma+1,2-\gamma,z)$，它和(d.1)式线性无关，故方程(d.2)的通解为
\begin{equation}
  u=C_1\mathrm{F}(\alpha,\gamma,z)
  +C_2z^{1-\gamma}\mathrm{F}(\alpha-\gamma+1,2-\gamma,z).
\end{equation}
第二项和第一项相反，在 $z=0$ 处有一奇点.

方程(d.2)为拉普拉斯型，它的解可用围线积分来表达.根据一般方法作函数
\begin{equation*}
  \mathrm{P}(t)=\gamma t-\alpha,\quad
  \mathrm{Q}(t)=t(t-1),\quad
  \mathrm{Z}(t)=t^{\alpha-1}(t-1)^{\gamma-\alpha-1},
\end{equation*}
则
\begin{equation}
  u=\int\mathrm{e}^{tz}t^{\alpha-1}(t-1)^{\gamma-\alpha-1}\,\mathrm{d}t.
\end{equation}
所选的积分路线必须使 $V(t)=\mathrm{e}^{tz}t^{\alpha}(t-1)^{\gamma-\alpha}$ 函数走完该曲线后回到原值.对(d.3)方程应用同一方法，可得 $u$ 的另一围线积分：
\begin{equation*}
  u=z^{1-\gamma}\int\mathrm{e}^{tz}t^{\alpha-\gamma}(t-1)^{-\alpha}\,\mathrm{d}t.
\end{equation*}
作替换 $tz\to t$ 后，这个积分化成下列方便形式：
\begin{equation}
  u(z)=\int\mathrm{e}^t(t-z)^{-\alpha}t^{\alpha-\gamma}\,\mathrm{d}t,
\end{equation}
而 $V$ 函数则变成
\begin{equation*}
  V(t)=\mathrm{e}^tt^{\alpha-\gamma+1}(t-z)^{1-\alpha}.
\end{equation*}

(d.6)中的被积式一般讲来具有两个奇点，在 $t=z$ 和 $t=0$ 处. 我们取积分围线 $C$，它从无穷远处（$\mathrm{Re}\ t\to-\infty$）来沿正方向绕过这两个奇点后，再回到无穷远处去（附录图5）.因为 $V(t)$ 在围线两端等于零，所以这个围线满足所需条件.(d.6)式沿围线 $C$ 的积分在 $z=0$ 处无奇点；因此，它除了一个常数因子外，必须和 $\mathrm{F}(\alpha,\gamma,z)$ 函数相同，$\mathrm{F}(\alpha,\gamma,z)$ 在 $z=0$ 处是没有奇点的.当 $z=0$ 时，被积函数的两个奇点重合；根据伽马函数论中的一个熟知公式\footnote{可参考 Whittaker 和 Watson, Course of Modern Analysis, Cambridge, 1944, §12.22.}，
\begin{equation}
  \frac{1}{2\pi i}\int_C\mathrm{e}^tt^{-\gamma}\,\mathrm{d}t=\frac{1}{\Gamma(\gamma)}.
\end{equation}
由于 $\mathrm{F}(\alpha,\gamma,0)=1$，显然有
\begin{equation}
  \mathrm{F}(\alpha,\gamma,z)=\frac{\Gamma(\gamma)}{2\pi i}\int_C
  \mathrm{e}^t(t-z)^{-\alpha}t^{\alpha-\gamma}\,\mathrm{d}t.
\end{equation}

(d.5) 中的被积函数在 $t=0$ 及 $t=1$ 处具有奇点. 如果 $\mathrm{Re}(\gamma-\alpha)>0$，并且 $\alpha$ 不是正整数，则积分路线可取围线 $C'$，它从 $t=1$ 点出发，沿正方向绕 $t=0$ 点再回到 $t=1$ 点处（附录图6）；当 $\mathrm{Re}(\gamma-\alpha)>0$ 时，函数 $V(t)$ 绕此围线一周而回到原来的零值\footnote{如果 $\gamma$ 为正整数，则 $C'$ 可取同时绕 $t=0$ 及 $t=1$ 点的任一围线.}.由此定义的积分在 $z=0$ 处也没有奇点，它和 $\mathrm{F}(\alpha,\gamma,z)$ 的关系为
\begin{equation}
  \mathrm{F}(\alpha,\gamma,z)=-\frac{1}{2\pi i}
  \frac{\Gamma(1-\alpha)\Gamma(\gamma)}{\Gamma(\gamma-\alpha)}
  \oint_{C'}\mathrm{e}^{tz}(-t)^{\alpha-1}(1-t)^{\gamma-\alpha-1}\,\mathrm{d}t.
\end{equation}

积分式(d.8)，(d.9)应作下列说明.当 $\alpha$ 和 $\gamma$ 为非整数时，被积函数不是单值函数.它们在每一点所取之值，已经假定满足这样的条件，即复量的幂次所取的辐角具有最小的绝对值.

注意下列有用的公式：
\begin{equation}
  \mathrm{F}(\alpha,\gamma,z)=\mathrm{e}^z\mathrm{F}(\gamma-\alpha,\gamma,-z),
\end{equation}
如果在积分式(d.8)中作替换 $t\to t+z$，立即可得上式.

我们曾经指出，如果 $\alpha=-n$（$n$ 为正整数），则函数 $\mathrm{F}(\alpha,\gamma,z)$ 变成一个多项式.可以求得这种多项式的简洁表达式.在积分式(d.9)中作替换 $t\to 1-t/z$ 后，再应用柯西公式，即得
\begin{equation}
  \mathrm{F}(-n,\gamma,z)=\frac{1}{\gamma(\gamma+1)\cdots(\gamma+n-1)}
  z^{1-\gamma}\mathrm{e}^z\frac{\mathrm{d}^n}{\mathrm{d}z^n}
  \left(\mathrm{e}^{-z}z^{\gamma+n-1}\right).
\end{equation}
如果尚有 $\gamma=m$（$m$ 为正整数），我们得下列公式
\begin{equation}
  \mathrm{F}(-n,m,z)=\frac{(-1)^{m-1}}{m(m+1)\cdots(m+n-1)}
  \mathrm{e}^z\frac{\mathrm{d}^{m+n-1}}{\mathrm{d}z^{m+n-1}}
  \left(\mathrm{e}^{-z}z^n\right).
\end{equation}
(d.8)式中作替换 $t\to z-t$ 后，再对积分式采用柯西公式，即得上式.

多项式 $\mathrm{F}(-n,m,z)$（$0\leqslant m\leqslant n$，除开一个常数因子外），与下列广义拉盖尔多项式相符；
\begin{equation}
\begin{aligned}
  \mathrm{L}_n^m(z)&=(-1)^m\frac{(n!)^2}{m!\,(n-m)!}
  \mathrm{F}\left(-(n-m),m+1,z\right)=\\
  &=\frac{n!}{(n-m)!}\mathrm{e}^z\frac{\mathrm{d}^n}{\mathrm{d}z^n}\mathrm{e}^{-z}z^{n-m}
  =(-1)^m\frac{n!}{(n-m)!}\mathrm{e}^z z^{-m}
  \frac{\mathrm{d}^{n-m}}{\mathrm{d}z^{n-m}}\mathrm{e}^{-z}z^n.
\end{aligned}
\end{equation}
$m=0$ 时 $\mathrm{L}_n^m$ 多项式用 $\mathrm{L}_n(z)$ 标记并称为拉盖尔多项式；由(d.13)式得
\begin{equation*}
  \mathrm{L}_n(z)=\mathrm{e}^z\frac{\mathrm{d}^n}{\mathrm{d}z^n}
  \left(\mathrm{e}^{-z}z^n\right).
\end{equation*}

(d.8)的积分表式便于得出 $z$ 很大时的合流超几何函数渐近式.我们把积分围线变形成 $C_1$ 和 $C_2$ 两条围线（图5），分别绕过 $t=0$ 和 $t=z$ 点；$C_2$ 的下支假定与 $C_1$ 的上支在无穷远处相接.我们在被积函数的括号内取出 $(-z)^{-\alpha}$，可得按 $z$ 的负幂次展开的表式.在沿围线 $C_2$ 的积分式中，作替换 $t\to t+z$；则围线 $C_2$ 变换成 $C_1$. 故(d.8)式可表成
\begin{equation}
\begin{aligned}
  \mathrm{F}(\alpha,\gamma,z)={}&\frac{\Gamma(\gamma)}{\Gamma(\gamma-\alpha)}
  (-z)^{-\alpha}\mathrm{G}(\alpha,\alpha-\gamma+1,-z)+\\
  &+\frac{\Gamma(\gamma)}{\Gamma(\alpha)}\mathrm{e}^z z^{\alpha-\gamma}
  \mathrm{G}(\gamma-\alpha,1-\alpha,z),
\end{aligned}
\end{equation}
其中的
\begin{equation}
  \mathrm{G}(\alpha,\beta,z)=\frac{\Gamma(1-\beta)}{2\pi i}\int_{c_1}
  \left(1+\frac{t}{z}\right)^{-\alpha}t^{\beta-1}\mathrm{e}^t\,\mathrm{d}t.
\end{equation}
(d.14)式中 $-z$ 和 $z$ 的幂次所取的辐角必须具有最小的绝对值.最后，被积函数中的 $\left(1+t/z\right)^{-\alpha}$ 展成 $t/z$ 的幂级数，并应用(d.7)式，可得 $\mathrm{G}(\alpha,\beta,z)$ 的渐近级数
\begin{equation}
  \mathrm{G}(\alpha,\beta,z)=1+\frac{\alpha\beta}{1!\,z}
  +\frac{\alpha(\alpha+1)\beta(\beta+1)}{2!\,z^2}+\cdots.
\end{equation}
(d.14)和(d.16)式给出了 $\mathrm{F}(\alpha,\gamma,z)$ 函数的渐近展式.

$\gamma$ 为正整数时，方程式(d.2)的通解(d.4)中的第二项，或则和第一项相符（如果 $\gamma=1$），或则毫无意义（如果 $\gamma>1$）. 在这种情形下，(d.14)式中的两项，即分别沿围线 $C_1$ 和 $C_2$ 的两个(d.8)积分式，可以取作(d.2)式的两个线性独立解［$C_1,C_2$ 围线和 $C$ 一样，满足所需条件，因此沿这两个围线的积分都是(d.2)式之解］.这两个解的渐近式可从早已得到的公式给出；留下的问题是求它们的 $z$ 的升幂展式.为此，我们从(d.14)式以及 $z^{1-\gamma}\mathrm{F}(\alpha-\gamma+1,2-\gamma,z)$ 所满足的类似公式出发.根据这两个公式，我们把 $\mathrm{G}(\alpha,\alpha-\gamma+1,-z)$ 表成 $\mathrm{F}(\alpha,\gamma,z)$ 和 $\mathrm{F}(\alpha-\gamma+1,2-\gamma,z)$；然后令 $\gamma=p+\varepsilon$（$p$ 是一个正整数），并取极限 $\varepsilon\to 0$，利用洛毕达定则求出此未定式.经过较长的计算后，给出下列展式：
\begin{equation}
\begin{aligned}
  \mathrm{G}(\alpha,\alpha-p+1,-z)={}&-\frac{\sin\pi\alpha\cdot\Gamma(p-\alpha)}{\pi\Gamma(p)}
  z^{\alpha}\Bigl\{\ln z\cdot\mathrm{F}(\alpha,p,z)+\\
  &+\sum_{s=0}^{\infty}
  \frac{\Gamma(p)\Gamma(\alpha+s)\left[\psi(\alpha+s)-\psi(p+s)-\psi(s+1)\right]}
  {\Gamma(\alpha)\Gamma(s+p)\Gamma(s+1)}z^s+\\
  &+\sum_{s=1}^{p-1}(-1)^{s+1}\frac{\Gamma(s)\Gamma(\alpha-s)\Gamma(p)}
  {\Gamma(\alpha)\Gamma(p-s)}z^{-s}\Bigr\},
\end{aligned}
\end{equation}
式中 $\psi$ 表示伽马函数的对数导数；$\psi(\alpha)=\Gamma'(\alpha)/\Gamma(\alpha)$.

\section{超几何函数}

超几何函数由 $|z|<1$ 圆内的下列级数定义：
\begin{equation}
  \mathrm{F}(\alpha,\beta,\gamma,z)=1+\frac{\alpha\beta}{\gamma}\frac{z}{1!}
  +\frac{\alpha(\alpha+1)\beta(\beta+1)}{\gamma(\gamma+1)}\frac{z^2}{2!}+\cdots,
\end{equation}
$|z|>1$ 时，由以上级数解析延拓求得［见(e.6)］.超几何函数是下列微分方程的一个特殊积分：
\begin{equation}
  z(1-z)u''+\left[\gamma-(\alpha+\beta+1)z\right]u'-\alpha\beta u=0.
\end{equation}
参量 $\alpha$ 和 $\beta$ 是任意的，而 $\gamma\neq 0,-1,-2,\cdots$，$\mathrm{F}(\alpha,\beta,\gamma,z)$ 函数显然对称于参量 $\alpha$ 和 $\beta$\footnote{合流超几何函数可由 $\mathrm{F}(\alpha,\beta,\gamma,z)$ 取下列极限求得
\begin{equation*}
  \mathrm{F}(\alpha,\gamma,z)=\lim_{\beta\to\infty}\mathrm{F}(\alpha,\beta,\gamma,z/\beta).
\end{equation*}
文献中还把超几何函数记作 $_2\mathrm{F}_1(\alpha,\beta,\gamma,z)$，合流超几何函数记作 $_1\mathrm{F}_1(\alpha,\gamma,z)$.F 下面的左标和右标分别代表级数项中分子上和分母上的参量数.}.(e.2)方程的第二个独立解为
\begin{equation*}
  z^{1-\gamma}\mathrm{F}(\beta-\gamma+1,\alpha-\gamma+1,2-\gamma,z);
\end{equation*}
它在 $z=0$ 处有一奇点.

为参考计，下面将给出超几何函数所满足的一些关系式.

$\mathrm{F}(\alpha,\beta,\gamma,z)$ 函数可以对所有的 $z$ 值表出，当 $\mathrm{Re}(\gamma-\alpha)>0$ 时，可表为下列积分式
\begin{equation}
  \mathrm{F}(\alpha,\beta,\gamma,z)=-\frac{1}{2\pi i}
  \frac{\Gamma(1-\alpha)\Gamma(\gamma)}{\Gamma(\gamma-\alpha)}
  \oint_{C'}(-t)^{\alpha-1}(1-t)^{\gamma-\alpha-1}(1-tz)^{-\beta}\,\mathrm{d}t,
\end{equation}
所取的围线 $C'$ 如附录图6所示.应用直接代入法容易验证这个积分式确实满足方程(e.2)；上式中所选的常数因子，应使得 $z=0$ 时函数值等于1.

方程(e.2)中作替换 $u=(1-z)^{\gamma-\alpha-\beta}u_1$，得到一个同一类型的方程，原式中的 $\alpha,\beta,\gamma$ 分别被 $\gamma-\alpha,\gamma-\beta,\gamma$ 诸参量所代替.因此，有等式
\begin{equation}
  \mathrm{F}(\alpha,\beta,\gamma,z)=(1-z)^{\gamma-\alpha-\beta}
  \mathrm{F}(\gamma-\alpha,\gamma-\beta,\gamma,z)
\end{equation}
上式两边满足同一方程式，并且在 $z=0$ 处它们的值相等.

将积分式(e.3)作替换 $t\to t/(1-z+zt)$，可得出以 $z$ 和 $z/(z-1)$ 为变量的两个超几何函数之间的下列关系式：
\begin{equation}
  \mathrm{F}(\alpha,\beta,\gamma,z)=(1-z)^{-\alpha}
  \mathrm{F}\left(\alpha,\gamma-\beta,\gamma,\frac{z}{z-1}\right).
\end{equation}
此式中的多值函数 $(1-z)^{-\alpha}$（及以后所有各式中类似的函数）所取之值，决定于这样的条件，即复量的幂所取的辐角具有最小的绝对值.

其次，我们给出以 $z$ 和 $1/z$ 为变量的超几何函数之间的下列重要公式，但不加证明.
\begin{equation}
\begin{aligned}
  \mathrm{F}(\alpha,\beta,\gamma,z)={}&\frac{\Gamma(\gamma)\Gamma(\beta-\alpha)}
  {\Gamma(\beta)\Gamma(\gamma-\alpha)}(-z)^{-\alpha}
  \mathrm{F}\left(\alpha,\alpha+1-\gamma,\alpha+1-\beta,\frac{1}{z}\right)+\\
  &+\frac{\Gamma(\gamma)\Gamma(\alpha-\beta)}
  {\Gamma(\alpha)\Gamma(\gamma-\beta)}(-z)^{-\beta}
  \mathrm{F}\left(\beta,\beta+1-\gamma,\beta+1-\alpha,\frac{1}{z}\right).
\end{aligned}
\end{equation}
这个公式把 $\mathrm{F}(\alpha,\beta,\gamma,z)$ 表为 $|z|>1$ 时为收敛的一个级数，这就是原级数(e.1)的解析延拓.

下列公式
\begin{equation}
\begin{aligned}
  \mathrm{F}(\alpha,\beta,\gamma,z)={}&\frac{\Gamma(\gamma)\Gamma(\gamma-\alpha-\beta)}
  {\Gamma(\gamma-\alpha)\Gamma(\gamma-\beta)}
  \mathrm{F}(\alpha,\beta,\alpha+\beta+1-\gamma,1-z)+\\
  &+\frac{\Gamma(\gamma)\Gamma(\alpha+\beta-\gamma)}
  {\Gamma(\alpha)\Gamma(\beta)}(1-z)^{\gamma-\alpha-\beta}
  \mathrm{F}(\gamma-\alpha,\gamma-\beta,\gamma+1-\alpha-\beta,1-z)
\end{aligned}
\end{equation}
把 $z$ 和 $1-z$ 的超几何函数联系起来；其推导过程与(e.6)式相似.把(e.7)和(e.6)式合并，可得下列关系式
\begin{equation}
\begin{aligned}
  \mathrm{F}(\alpha,\beta,\gamma,z)={}&\frac{\Gamma(\gamma)\Gamma(\beta-\alpha)}
  {\Gamma(\beta)\Gamma(\gamma-\alpha)}(1-z)^{-\alpha}
  \mathrm{F}\left(\alpha,\gamma-\beta,\alpha+1-\beta,\frac{1}{1-z}\right)+\\
  &+\frac{\Gamma(\gamma)\Gamma(\alpha-\beta)}
  {\Gamma(\alpha)\Gamma(\gamma-\beta)}(1-z)^{-\beta}
  \mathrm{F}\left(\beta,\gamma-\alpha,\beta+1-\alpha,\frac{1}{1-z}\right),
\end{aligned}
\end{equation}
\begin{equation}
\begin{aligned}
  \mathrm{F}(\alpha,\beta,\gamma,z)={}&\frac{\Gamma(\gamma)\Gamma(\gamma-\alpha-\beta)}
  {\Gamma(\gamma-\beta)\Gamma(\gamma-\alpha)}z^{-\alpha}
  \mathrm{F}\left(\alpha,\alpha+1-\gamma,\alpha+\beta+1-\gamma,\frac{z-1}{z}\right)+\\
  &+\frac{\Gamma(\gamma)\Gamma(\alpha+\beta-\gamma)}
  {\Gamma(\alpha)\Gamma(\beta)}(1-z)^{\gamma-\alpha-\beta}z^{\alpha-\gamma}
  \mathrm{F}\left(1-\beta,\gamma-\beta,\gamma+1-\alpha-\beta,\frac{z-1}{z}\right).
\end{aligned}
\end{equation}
(e.6)到(e.9)各式右边的每一项都是超几何方程的一个解.

如果 $\alpha$（或 $\beta$）是一个负整数（或零），$\alpha=-n$，则超几何函数变成 $n$ 次多项式，并可表成下列形式
\begin{equation}
  \mathrm{F}(-n,\beta,\gamma,z)=\frac{z^{1-\gamma}(1-z)^{\gamma+n-\beta}}
  {\gamma(\gamma+1)\cdots(\gamma+n-1)}
  \frac{\mathrm{d}^n}{\mathrm{d}z^n}
  \left[z^{\gamma+n-1}(1-z)^{\beta-\gamma}\right].
\end{equation}
此式与下列雅可比多项式只差一个常数因子.
\begin{equation}
\begin{aligned}
  \mathrm{P}_n^{(a,b)}(z)&=\frac{(a+1)(a+2)\cdots(a+n)}{n!}
  \mathrm{F}\left(-n,a+b+n+1,a+1,\frac{1-z}{2}\right)=\\
  &=\frac{(-1)^n}{2^n n!}(1-z)^{-a}(1+z)^{-b}\frac{\mathrm{d}^n}{\mathrm{d}z^n}
  \left[(1-z)^{a+n}(1+z)^{b+n}\right].
\end{aligned}
\end{equation}
$a=b=0$ 时，雅可比多项式就是勒让德多项式. $n=0$ 时，$\mathrm{P}_0^{(a,b)}=1$.

\section{含有合流超几何函数的积分计算}

让我们考虑下列积分式
\begin{equation}
  \mathrm{J}_{\alpha\gamma}^{\nu}=\int_0^{\infty}\mathrm{e}^{-\lambda z}z^{\nu}
  \mathrm{F}(\alpha,\gamma,kz)\,\mathrm{d}z,
\end{equation}
假定它是收敛的，那么，必须使 $\mathrm{Re}\ \nu>-1$ 和 $\mathrm{Re}\ \lambda>|\mathrm{Re}\ k|$；如果 $\alpha$ 是一个负整数，后一条件可用 $\mathrm{Re}\ \lambda>0$ 代替.

应用 $\mathrm{F}(\alpha,\gamma,kz)$ 的积分表式(d.9)，并把围线积分和与 $\mathrm{d}z$ 的积分对调，容易得出：
\begin{equation*}
\begin{aligned}
  \mathrm{J}_{\alpha\gamma}^{\nu}={}&-\frac{1}{2\pi i}
  \frac{\Gamma(1-\alpha)\Gamma(\gamma)}{\Gamma(\gamma-\alpha)}
  \oint_{C'}\int_0^{\infty}\mathrm{e}^{-(\lambda-kt)z}z^{\nu}
  (-t)^{\alpha-1}(1-t)^{\gamma-\alpha-1}\,\mathrm{d}t\,\mathrm{d}z=\\
  ={}&-\frac{1}{2\pi i}\frac{\Gamma(1-\alpha)\Gamma(\gamma)}{\Gamma(\gamma-\alpha)}
  \lambda^{-\nu-1}\Gamma(\nu+1)\oint_{C'}(-t)^{\alpha-1}(1-t)^{\gamma-\alpha-1}
  \left(1-\frac{kt}{\lambda}\right)^{-\nu-1}\mathrm{d}t.
\end{aligned}
\end{equation*}
考虑到(e.3)式，我们最后得
\begin{equation}
  \mathrm{J}_{\alpha\gamma}^{\nu}=\Gamma(\nu+1)\lambda^{-\nu-1}
  \mathrm{F}(\alpha,\nu+1,\gamma,k/\lambda)
\end{equation}
当 $\mathrm{F}(\alpha,\nu+1,\gamma,k/\lambda)$ 函数变成一个多项式时，积分 $\mathrm{J}_{\alpha\gamma}^{\nu}$ 可用初等函数表示出来：
\begin{equation}
  \mathrm{J}_{\alpha\gamma}^{\nu+n-1}=(-1)^n\Gamma(\gamma)
  \frac{\mathrm{d}^n}{\mathrm{d}\lambda^n}
  \left[\lambda^{\alpha-\gamma}(\lambda-k)^{-\alpha}\right],
\end{equation}
\begin{equation}
  \mathrm{J}_{-n,\gamma}^{\nu}=(-1)^n
  \frac{\Gamma(\nu+1)(\lambda-k)^{\gamma+n-\nu-1}}
  {\gamma(\gamma+1)\cdots(\gamma+n-1)}
  \frac{\mathrm{d}^n}{\mathrm{d}\lambda^n}
  \left[\lambda^{-\nu-1}(\lambda-k)^{\nu-\gamma+1}\right],
\end{equation}
\begin{equation}
\begin{aligned}
  \mathrm{J}_{\alpha m}^{n}={}&\frac{(-1)^{m-n}}
  {k^{m-1}(1-\alpha)(2-\alpha)\cdots(m-1-\alpha)}\times\\
  &\times\Bigl\{-(m-1)!\frac{\mathrm{d}^{n}}{\mathrm{d}\lambda^{n}}
  \left[\lambda^{\alpha-1}(\lambda-k)^{m-\alpha-1}\right]+\\
  &+n!\,(m-n-1)\cdots(m-1)\lambda^{\alpha-n-1}(\lambda-k)^{-1+m-n-\alpha}\times\\
  &\times\frac{\mathrm{d}^{m-n-2}}{\mathrm{d}\lambda^{m-n-2}}
  \left[\lambda^{m-\alpha-1}(\lambda-k)^{\alpha-1}\right]\Bigr\}
\end{aligned}
\end{equation}
其中 $m,n$ 都是整数，且 $0\leqslant n\leqslant m-2$.

其次，让我们计算下列积分
\begin{equation}
  \mathrm{J}_{\nu}=\int_0^{\infty}\mathrm{e}^{-kz}z^{\nu-1}
  \left[\mathrm{F}(-n,\gamma,kz)\right]^2\mathrm{d}z,
\end{equation}
其中 $n$ 是一个正整数，并且 $\mathrm{Re}\ \nu>0$.为了计算这个积分式，我们从一个更普遍的积分式出发，被积函数中的 $\mathrm{e}^{-kz}$ 因子改成 $\mathrm{e}^{-\lambda z}$.其中一个 $\mathrm{F}(-n,\gamma,kz)$ 函数可写成围线积分(d.9)，然后应用(f.3)式对 $\mathrm{d}z$ 积分：
\begin{equation*}
\begin{aligned}
  &\int_0^{\infty}\mathrm{e}^{-\lambda z}z^{\nu-1}
  \left[\mathrm{F}(-n,\gamma,kz)\right]^2\mathrm{d}z
  =-\frac{1}{2\pi i}\frac{\Gamma(1+n)\Gamma(\gamma)}{\Gamma(\gamma+n)}\times\\
  &\qquad\times\int_0^{\infty}\oint_{C'}(-t)^{-n-1}(1-t)^{\gamma+n-1}
  \mathrm{e}^{-(\lambda-kt)z}z^{\nu-1}\mathrm{F}(-n,\gamma,kz)\,\mathrm{d}t\,\mathrm{d}z=\\
  &\qquad=-\frac{1}{2\pi i}(-1)^n
  \frac{\Gamma(1+n)\Gamma^2(\gamma)\Gamma(\nu)}{\Gamma^2(\gamma+n)}\times\\
  &\qquad\qquad\times\oint_{C'}(\lambda-kt-k)^{\gamma+n-\nu}(-t)^{-n-1}\times\\
  &\qquad\qquad\times(1-t)^{\gamma+n-1}
  \frac{\mathrm{d}^n}{\mathrm{d}\lambda^n}
  \left[(\lambda-kt)^{-\nu}(\lambda-kt-k)^{\nu-\gamma}\right]\mathrm{d}t.
\end{aligned}
\end{equation*}
$\lambda$ 的 $n$ 次导数显然可用 $t$ 的 $n$ 次导数来代替；然后再令 $\lambda=k$，就回到积分 $\mathrm{J}_{\nu}$
\begin{equation*}
\begin{aligned}
  \mathrm{J}_{\nu}={}&-\frac{1}{2\pi i}
  \frac{\Gamma(n+1)\Gamma(\nu)\Gamma^2(\gamma)}{\Gamma^2(\gamma+n)k^{\nu}}\times\\
  &\times\oint_{C'}(-t)^{\gamma-\nu-1}(1-t)^{\gamma+n-1}
  \frac{\mathrm{d}^n}{\mathrm{d}t^n}
  \left[(1-t)^{-\nu}(-t)^{\nu-\gamma}\right]\mathrm{d}t.
\end{aligned}
\end{equation*}
分部积分 $n$ 次后. $\mathrm{d}^n/\mathrm{d}t^n$ 算符就转移到表式 $(-t)^{\gamma-\nu-1}(1-t)^{\gamma+n-1}$ 身上，然后用莱布尼茨公式展开此导数.结果得到许多积分式之和，其中的每一个积分式都能化成熟知的欧拉积分式.最后可得所求积分的下列表式：
\begin{equation}
\begin{aligned}
  \mathrm{J}_{\nu}={}&\frac{\Gamma(\nu)\,n!}
  {k^{\nu}\gamma(\gamma+1)\cdots(\gamma+n-1)}\times\\
  &\times\Bigl\{1+\sum_{s=0}^{n-1}
  \frac{n(n-1)\cdots(n-s)(\gamma-\nu-s-1)(\gamma-\nu-s)\cdots(\gamma-\nu+s)}
  {\left[(s+1)!\right]^2\gamma(\gamma+1)\cdots(\gamma+s)}\Bigr\}
\end{aligned}
\end{equation}
容易看出 $\mathrm{J}_{\nu}$ 积分间具有下列关系式
\begin{equation}
  \mathrm{J}_{\gamma+p}=\frac{(\gamma-p-1)(\gamma-p)\cdots(\gamma+p-1)}{k^{2p+1}}
  \mathrm{J}_{\gamma-1-p},
\end{equation}
其中 $p$ 是任意整数.

同样地可计算积分
\begin{equation}
  \mathrm{J}=\int_0^{\infty}\mathrm{e}^{-\lambda z}z^{\gamma-1}
  \mathrm{F}(\alpha,\gamma,kz)\mathrm{F}(\alpha',\gamma,k'z)\,\mathrm{d}z,
\end{equation}
我们把 $\mathrm{F}(\alpha',\gamma,k'z)$ 函数表成围线积分(d.9)，然后利用 $n=0$ 的(f.3)式对 $\mathrm{d}z$ 积分：
\begin{equation*}
\begin{aligned}
  \mathrm{J}={}&-\frac{1}{2\pi i}\frac{\Gamma(1-\alpha')\Gamma(\gamma)}
  {\Gamma(\gamma-\alpha')}
  \oint_{C'}\int_0^{\infty}(-t)^{\alpha'-1}(1-t)^{\gamma-\alpha'-1}z^{\gamma-1}
  \mathrm{e}^{-z(\lambda-k')}\times\\
  &\times\mathrm{F}(\alpha,\gamma,kz)\,\mathrm{d}z\,\mathrm{d}t\\
  ={}&-\frac{1}{2\pi i}\frac{\Gamma(1-\alpha')\Gamma^2(\gamma)}
  {\Gamma(\gamma-\alpha')}
  \oint_{C'}(-t)^{\alpha'-1}(1-t)^{\gamma-\alpha'-1}
  (\lambda-k't)^{\alpha-\gamma}\times\\
  &\times(\lambda-k't-k)^{-\alpha}\,\mathrm{d}t.
\end{aligned}
\end{equation*}
作替换 $t\to\lambda t/(k't+\lambda-k')$，这个积分变成(e.3)形式，结果得
\begin{equation}
  \mathrm{J}=\Gamma(\gamma)\lambda^{\alpha+\alpha'-\gamma}(\lambda-k)^{-\alpha}
  (\lambda-k')^{-\alpha'}
  \mathrm{F}\left(\alpha,\alpha',\gamma,
  \frac{kk'}{(\lambda-k)(\lambda-k')}\right).
\end{equation}
如果 $\alpha$（或 $\alpha'$）是负整数，$\alpha=-n$，上式可用(e.7)式改写成
\begin{equation}
\begin{aligned}
  \mathrm{J}={}&\frac{\Gamma^2(\gamma)\Gamma(\gamma+n-\alpha')}
  {\Gamma(\gamma+n)\Gamma(\gamma-\alpha')}
  \lambda^{-n+\alpha'-\gamma}(\lambda-k)^n(\lambda-k')^{-\alpha'}\times\\
  &\times\mathrm{F}\left(-n,\alpha',-n+\alpha'+1-\gamma,
  \frac{\lambda(\lambda-k-k')}{(\lambda-k)(\lambda-k')}\right).
\end{aligned}
\end{equation}

最后，让我们考虑下列形式的积分
\begin{equation}
  \mathrm{J}_{\gamma}^{sp}(\alpha,\alpha')=\int_0^{\infty}
  \mathrm{e}^{-\frac{1}{2}(k+k')z}z^{\gamma-1+s}
  \mathrm{F}(\alpha,\gamma,kz)\mathrm{F}(\alpha',\gamma-p,k'z)\,\mathrm{d}z.
\end{equation}
假定参量值使该积分绝对收敛；$s$ 和 $p$ 都是正整数.应用(f.10)，其中最简单的积分式 $\mathrm{J}_{\gamma}^{00}(\alpha,\alpha')$ 为
\begin{equation}
\begin{aligned}
  \mathrm{J}_{\gamma}^{00}(\alpha,\alpha')={}&2^{\gamma}\Gamma(\gamma)
  (k+k')^{\alpha+\alpha'-\gamma}(k'-k)^{-\alpha}(k-k')^{-\alpha'}\times\\
  &\times\mathrm{F}\left[\alpha,\alpha',\gamma,\frac{-4kk'}{(k'-k)^2}\right],
\end{aligned}
\end{equation}
如果 $\alpha$（或 $\alpha'$）是一个负整数，$\alpha=-n$，应用(f.11)式，我们还可写成
\begin{equation}
\begin{aligned}
  \mathrm{J}_{\gamma}^{00}(-n,\alpha')={}&2^{\gamma}
  \frac{\Gamma(\gamma)(\gamma-\alpha')(\gamma-\alpha'+1)\cdots(\gamma-\alpha'+n-1)}
  {\gamma(\gamma+1)\cdots(\gamma+n-1)}\times\\
  &\times(-1)^n(k+k')^{-n+\alpha'-\gamma}(k-k')^{n-\alpha'}
  \mathrm{F}\left[-n,\alpha',\alpha'+1-n-\gamma,
  \left(\frac{k+k'}{k-k'}\right)^2\right].
\end{aligned}
\end{equation}
$\mathrm{J}_{\gamma}^{sp}(\alpha,\alpha')$ 的一般公式也可以推导出来，但是它过于复杂不便应用.更方便的办法是应用递推公式，它可以把 $\mathrm{J}_{\gamma}^{sp}(\alpha,\alpha')$ 的积分式化成 $s=p=0$ 的积分式. 我们把它写出来，不加证明\footnote{具体推导见 W.~Gordon, Annalen der Physik [5]2, 1031, 1929.}.公式
\begin{equation}
  \mathrm{J}_{\gamma}^{sp}(\alpha,\alpha')=\frac{\gamma-1}{k}
  \left\{\mathrm{J}_{\gamma-1}^{s,p-1}(\alpha,\alpha')
  -\mathrm{J}_{\gamma-1}^{s,p-1}(\alpha-1,\alpha')\right\},
\end{equation}
可使 $\mathrm{J}_{\gamma}^{sp}(\alpha,\alpha')$ 化成 $p=0$ 的积分式.然后，公式
\begin{equation}
\begin{aligned}
  \mathrm{J}_{\gamma}^{s+1,0}(\alpha,\alpha')=\frac{1}{k^2-k'^2}\Bigl\{
  &\left[\frac{\gamma}{2}(k-k')-k\alpha+k'\alpha'-k's\right]
  \mathrm{J}_{\gamma}^{s0}(\alpha,\alpha')+\\
  &+s(\gamma-1+s-2\alpha')\mathrm{J}_{\gamma}^{s-1,0}(\alpha,\alpha')
  +2\alpha's\,\mathrm{J}_{\gamma}^{s-1,0}(\alpha,\alpha'+1)\Bigr\}
\end{aligned}
\end{equation}
最后可使 $\mathrm{J}_{\gamma}^{s0}(\alpha,\alpha')$ 化成 $s=p=0$ 的积分式.

% （附录为全书最后一个文件，其后无正文；为稳妥起见在此恢复
%   main.tex 中的节号与公式编号定义）
\ctexset{section={number=\arabic{section}}}
\renewcommand{\thesection}{\arabic{section}}
\renewcommand{\theequation}{\arabic{section}.\arabic{equation}}
